Dependable robots in industrial settings must do more than keep a policy and controller internally consistent: they must operate within task procedures, hardware limits, and constraints created by shared human workspaces. My representation research and CR5 deployment give me a technical starting point, but not yet the human-aware and operational perspective required for such systems. The MSc in Intelligent Robotics Engineering at The Hong Kong Polytechnic University (PolyU) would let me develop that perspective across learning, control, and robotic operation.

The published 2027/28 curriculum lists Embodied Robot Intelligence and Advanced Control Technology as core subjects, which would connect learned actions to physical behavior and deepen the feedback reasoning behind my ServoP decision. If offered when I register, the elective Industrial Human-Robot Systems and Automation would add human-aware workflows and industrial operating constraints that my previous projects did not address. If I were eligible for and selected the dissertation route, its seven taught subjects plus a Dissertation could support a study of how learned robot behavior should adapt when operating procedures, workspace limits, and task allocation change. I would contribute experience tracing behavior across synchronized data, policy interfaces, and robot commands, while using the programme to learn how those technical layers should adapt in human-aware industrial operation.
